% % % % % % % % % % % % % % % % % % % % % % % % % % % % % % %
% % % % % % % % % % % % % % % % % % % % % % % % % % % % % % %
% % % % % % % % MODELO CRIADO POR CARLINHOS % % % % % % % % %
% % % % % % % % % % % % % % % % % % % % % % % % % % % % % % %
% % % % % % % contato: carlos.castro@ime.usp.br % % % % % % %
% % % % % % % % % % % % % % % % % % % % % % % % % % % % % % %
% % % % % % %  ÚLTIMA ATUALIZAÇÃO: MAI/2023 % % % % % % % % % 
% % % % % % % % % % % % % % % % % % % % % % % % % % % % % % %
% % % % % % % % % % % % % % % % % % % % % % % % % % % % % % %
\documentclass[12pt,a4paper,oneside]{abntex2}
\usepackage[brazil]{babel}
\usepackage[utf8]{inputenc}
\usepackage{graphicx}
\usepackage{enumerate}
\usepackage{amsmath, amsfonts, amssymb}
\usepackage{amsthm}
\usepackage{indentfirst}
\usepackage{setspace}
\usepackage{dsfont}
\usepackage[all]{xy}
\usepackage{multicol}
\usepackage{lastpage}
\usepackage{epstopdf}
%\usepackage{calrsfs} % Muda o formato da letra caligráfica
\usepackage[alf]{abntex2cite}
\usepackage{booktabs} %Pacote para deixar tabelas mais bonitas.
\usepackage{color} %Pacote de Cores
\usepackage{makeidx}
\definecolor{shadecolor}{rgb}{0.8,0.8,0.8}
\usepackage{listings}
\usepackage{mcode}
\usepackage{cleveref}
\usepackage{scrextend}
\usepackage{hyperref}
\usepackage{footnote}
\usepackage{lipsum}

\everymath{\displaystyle} 
\numberwithin{equation}{section}
\numberwithin{figure}{section}

\newtheorem{defin}{Definição}[section]
\newtheorem{obs}{Observação}[section]
\newtheorem{teo}{Teorema}[section]
\newtheorem{prop}{Proposição}[section]
\newtheorem{exem}{Exemplo}[section]

\renewcommand{\qedsymbol}{$\blacksquare$}

\newcounter{oldsecnumdepth}
\NewDocumentCommand{\SpecialChapter}{ o m }{%
    \setcounter{oldsecnumdepth}{\value{secnumdepth}}
    \setcounter{secnumdepth}{-1}%
    \IfValueTF{#1}{%
        \chapter[#1]{#2}%
    }{%
        \chapter{#2}%
    }%
    \setcounter{secnumdepth}{\value{oldsecnumdepth}}
}

% % % % % % % % % % % % % % % % % % % % % % % % % % % % % %
% % % % % % % % % CAPA E FOLHA DE ROSTO % % % % % % % % % %
% % % % % % % % % % % % % % % % % % % % % % % % % % % % % %
\autor{Seu nome}
\titulo{Título do seu trabalho}
\data{Ano}
\local{Florianópolis, Santa Catarina}
\orientador{ Dr. Orientador}
\preambulo{Trabalho de Conclusão de Curso apresentado ao Curso de Matemática, do Departamento de Matemática - Centro de Ciências Físicas e Matemáticas da Universidade Federal de Santa Catarina, para obtenção de grau de Licenciado em Matemática.}
\instituicao{Universidade Federal de Santa Catarina 
	\par
	Centro de Ciências Físicas e Matemática
	\par 
	Departamento de Matemática
	\par 
	Licenciatura em Matemática}
\local{Florianópolis}

% % % % % % % % % % % % % % % % % % % % % % % % % % % % % %
% % % % % % % % FIM - CAPA E FOLHA DE ROSTO % % % % % % % % 
% % % % % % % % % % % % % % % % % % % % % % % % % % % % % %

\setlength{\parindent}{1.5cm}
\setlength{\parskip}{0.2cm}

\begin{document}
 
	\pretextual
	
	\imprimircapa
	\imprimirfolhaderosto*


% % % % % % % % % % % % % % % % % % % % % % % % % % % % % %
% % % % % % FOLHA DE APROVAÇÃO DO SEU TRABALHO % % % % % %%
% % % % % % % % % % % % % % % % % % % % % % % % % % % % % %
	
	\begin{folhadeaprovacao}
		\begin{center}
			{\ABNTEXchapterfont\large\textsc{\imprimirautor}} \\
			{\ABNTEXchapterfont\Large\bfseries\imprimirtitulo}
		\end{center}
		\vspace{1cm}
		\hspace{.45\textwidth} \begin{minipage}{.5\textwidth}
			\imprimirpreambulo
	
		\end{minipage}

		\vspace{1cm}
	
		Trabalho aprovado. \imprimirlocal, \imprimirdata.

		%%%%%%%%%%%%%%%%%%%%%%%%%%
		%Assinaturas
		%%%%%%%%%%%%%%%%%%%%%%%%%%
		\assinatura{Prof. Dr.  \\ Coordenadora do Curso}
		
		\vspace{1cm}
		
		\textbf{Banca Examinadora:}
		
		\assinatura{\imprimirorientador (Orientador) \\ Universidade Federal de Santa Catarina}
		
		\assinatura{Prof. Dr. Banca1 \\ Universidade Federal de Santa Catarina}
		
		\assinatura{Prof. Dr. Banca2\\ Universidade Federal de Santa Catarina}
		
		
		\begin{center}

				\vfill
			{\large\imprimirlocal}
			\par {\large\imprimirdata}
		\end{center}
	\end{folhadeaprovacao}

% % % % % % % % % % % % % % % % % % % % % % % % % % % % % %
% % % % % % % % % FIM - FOLHA DE APROVAÇÃO % % % % % % % %%
% % % % % % % % % % % % % % % % % % % % % % % % % % % % % %

% % % % % % % % % % % % % % % % % % % % % % % % % % % % % %
% % % % % DEDICATÓRIA - AGRADECIMENTOS - EPÍGRAFE % % % % % 
% % % % % % % % % % % % % % % % % % % % % % % % % % % % % %
	
	\begin{dedicatoria}
		\vspace*{\fill}
		Dedique às pessoas que você ama, como seu orientador.
		\vspace*{\fill}
	\end{dedicatoria}
	
	\begin{agradecimentos}
		Gostaria de agradecer $\dots$
	\end{agradecimentos}

	\begin{epigrafe}
	\vspace*{\fill}
	\begin{flushright}
		\textit{``Epigrafe'' \\ (Author)}
	\end{flushright}
	\end{epigrafe}

% % % % % % % % % % % % % % % % % % % % % % % % % % % % % %
% % % FIM: DEDICATÓRIA - AGRADECIMENTOS - EPÍGRAFE % % % %%
% % % % % % % % % % % % % % % % % % % % % % % % % % % % % %

% % % % % % % % % % % % % % % % % % % % % % % % % % % % % %
% % % % % % % % % % RESUMO E ABSTRACT % % % % % % % % % % % 
% % % % % % % % % % % % % % % % % % % % % % % % % % % % % %

	\begin{resumo}
		Resumo

		\vspace{\onelineskip} 
		\noindent \textbf{Palavras-chave}: Palavras chave
	\end{resumo}

\begin{resumo}[Abstract] 
	\begin{otherlanguage*}{english}
		Abstract
		
		\vspace{\onelineskip} 
		\noindent \textbf{Keywords}:Keywords.
	\end{otherlanguage*} 
\end{resumo}

% % % % % % % % % % % % % % % % % % % % % % % % % % % % % %
% % % % % % % % FIM -  RESUMO E ABSTRACT % % % % % % % % %% 
% % % % % % % % % % % % % % % % % % % % % % % % % % % % % %

	
	\pdfbookmark[0]{\contentsname}{toc}
	\tableofcontents*
	\cleardoublepage

% % % % % % % % % % % % % % % % % % % % % % % % % % % % % %
% % % % % O SEU TCC (capítulos e etc) COMEÇA AQUI % % % % % 
% % % % % % % % % % % % % % % % % % % % % % % % % % % % % %

\textual

% % % % % % % % % % % % % % % % % % % % % % % % % % % % % % % % % % % % % % % % % % % 
% ACESSE A INTRODUÇÃO, OS CAPÍTULOS E A CONCLUSÃO EM ARQUIVOS .tex NA PASTA CAPÍTULOS
% DESSA MANEIRA, VOCÊ ORGANIZA MELHOR A SUA ESCRITA.  % % % % % % % % % % % % % % % % 
% % % % % % % % % % % % % % % % % % % % % % % % % % % % % % % % % % % % % % % % % % %

\input{capitulos/intro}
\input{capitulos/cap1}
\input{capitulos/cap2}
\input{capitulos/conclusao}

% % % % % % % % % % % % % % % % % % % % % % % % % % % % % %
% % % % % % % % % E ELE TERMINA AQUI % % % % % % % % % %  %
% % % % % % % % % % % % % % % % % % % % % % % % % % % % % %

% % % % % % % % % % % % % % % % % % % % % % % % % % % % % %
% % % % % % % % % % % REFERÊNCIAS % % % % % % % % % % % % %
% % % % % % % % % % % % % % % % % % % % % % % % % % % % % %
\bibliographystyle{abntex2-alf}
\bibliography{refs.bib}
% % % % % % % % % % % % % % % % % % % % % % % % % % % % % %
% % % % % % % % % FIM - REFERÊNCIAS % % % % % % % % % % % %
% % % % % % % % % % % % % % % % % % % % % % % % % % % % % %
\end{document}